\newpage
\phantomsection
\addcontentsline{toc}{section}{Resumen}
\section*{Resumen}
La presente propuesta de plan de tesis aborda el diseño, implementación y evaluación de un Small Language Model bilingüe español-Quechua Collao para asistencia educativa offline en contextos rurales de Educación Intercultural Bilingüe del sur andino peruano. El problema se ubica en la convergencia de tres factores: la brecha de conectividad que afecta a más del 79\% de los hogares rurales sin acceso a internet, la limitada disponibilidad de recursos tecnológicos adaptados a lenguas indígenas de bajos recursos y la necesidad de apoyo pedagógico contextualizado para estudiantes que interactúan en español, Quechua Collao o formas mixtas de comunicación. La investigación plantea un asistente educativo que no funcione como traductor aislado, sino como un sistema capaz de detectar el idioma predominante de la consulta, recuperar contenidos curriculares y glosarios locales mediante generación aumentada por recuperación, y producir respuestas pedagógicas pertinentes en el idioma requerido. El estudio se delimita exclusivamente a la variedad Quechua Collao, de acuerdo con la ubicación geográfica del estudio, la disponibilidad de corpus paralelos validados y la presencia de rasgos fonológicos diferenciadores como consonantes oclusivas glotizadas y aspiradas. Se espera que la propuesta contribuya teórica, práctica y metodológicamente al integrar procesamiento de lenguaje natural para lenguas indígenas, modelos compactos desplegables en dispositivos locales mediante técnicas de cuantización y criterios pedagógicos propios de la Educación Intercultural Bilingüe.\\

\noindent
\textbf{Palabras clave}: Small Language Model, Quechua Collao, Educación Intercultural Bilingüe, asistencia educativa, inferencia offline, RAG.
